\exo
On considère les points $A(4;5)$ et $B(3;-1)$.
\begin{enumerate}
\item Déterminer un vecteur directeur de la droite $(AB)$.
\item Parmi les vecteurs suivants, lesquels sont des vecteurs directeurs de $(AB)$ ?

 $\overrightarrow{u} \left( \begin{array}{c}
2\\
12
\end{array} \right)$ $\overrightarrow{v} \left( \begin{array}{c}
-1\\
6
\end{array} \right)$ $\overrightarrow{w} \left( \begin{array}{c}
\dfrac{4}{3}\\
8
\end{array} \right)$ $\overrightarrow{z} \left( \begin{array}{c}
2,3\\
14
\end{array} \right)$
\end{enumerate}

\exo
On considère la droite $(AB)$ ci-dessous.


\definecolor{ududff}{rgb}{0.30196078431372547,0.30196078431372547,1}
\begin{tikzpicture}[scale=0.7,line cap=round,line join=round,>=triangle 45,x=1cm,y=1cm]
\begin{axis}[
x=1cm,y=1cm,
axis lines=middle,
ymajorgrids=true,
xmajorgrids=true,
xmin=-5.014398203376298,
xmax=9.664378229870339,
ymin=-2.3872017134381927,
ymax=5.628786354405039,
xtick={-5,-4,...,9},
ytick={-2,-1,...,5},]
\clip(-5.014398203376298,-2.3872017134381927) rectangle (9.664378229870339,5.628786354405039);
\draw [line width=0.8pt,domain=-5.014398203376298:9.664378229870339] plot(\x,{(--7--1*\x)/3});
\begin{scriptsize}
\draw [fill=ududff] (-1,2) circle (2.5pt);
\draw[color=ududff] (-0.9089279197026294,2.2515209953494857) node {$A$};
\draw [fill=ududff] (2,3) circle (2.5pt);
\draw[color=ududff] (2.0956341314775417,3.249219081046712) node {$B$};
\end{scriptsize}
\end{axis}
\end{tikzpicture}


\begin{enumerate}
\item Déterminer deux vecteurs directeurs de $(AB)$.
\item Parmi les vecteurs suivants, lesquels sont des vecteurs directeurs de $(AB)$ ?

 $\overrightarrow{u} \left( \begin{array}{c}
9\\
3
\end{array} \right)$ $\overrightarrow{v} \left( \begin{array}{c}
-11\\
-4
\end{array} \right)$ 
\end{enumerate}

\exo
On a tracé deux droites $\Delta$ et $\Delta'$. Donner un vecteur directeur de chacune d'elles.

\begin{tikzpicture}[scale=0.7,line cap=round,line join=round,>=triangle 45,x=1cm,y=1cm]
\begin{axis}[
x=1cm,y=1cm,
axis lines=middle,
ymajorgrids=true,
xmajorgrids=true,
xmin=-5.082954166586103,
xmax=7.856287069953209,
ymin=-3.2926945949362265,
ymax=3.7733441740801146,
xtick={-5,-4,...,7},
ytick={-3,-2,...,3},]
\clip(-5.082954166586103,-3.2926945949362265) rectangle (7.856287069953209,3.7733441740801146);
\draw [line width=2pt,domain=-5.082954166586103:7.856287069953209] plot(\x,{(--3--1*\x)/3});
\draw [line width=2pt,domain=-5.082954166586103:7.856287069953209] plot(\x,{(--2-4*\x)/2});
\begin{scriptsize}
\draw[color=black] (-4.365230629246813,-1.1445783740263948) node {$\Delta$};
\draw[color=black] (-0.30150017839618526,3.2830682366018706) node {$\Delta'$};
\end{scriptsize}
\end{axis}
\end{tikzpicture}

\begin{enumerate}
\item Pour chacun des vecteurs suivants, indiquer s'il est un vecteur directeur de $\Delta$, de $\Delta'$ ou d'aucune des deux droites.

 $\overrightarrow{u} \left( \begin{array}{c}
-\dfrac{3}{2}\\
3
\end{array} \right)$ $\overrightarrow{v} \left( \begin{array}{c}
-10\\
-2,3
\end{array} \right)$ $\overrightarrow{w} \left( \begin{array}{c}
87\\
29
\end{array} \right)$ $\overrightarrow{z} \left( \begin{array}{c}
\dfrac{1}{5}\\
-0,4
\end{array} \right)$

\end{enumerate}

\exo
\begin{enumerate}
\item Dans un repère, tracer la droite $d$ passant par $E(-1;4)$ et de vecteur directeur $\overrightarrow{v} \left( \begin{array}{c}
2\\
-3
\end{array} \right)$.
\item Tracer la droite $d'$ passant par $F(-3;0)$ et de vecteur directeur $\overrightarrow{u} \left( \begin{array}{c}
-3\\
\dfrac{9}{2}
\end{array} \right)$.
\item Démontrer que $\overrightarrow{u}$ et $\overrightarrow{v}$ sont colinéaires.
\item Que peut-on en déduire pour les droites $d$ et $d'$ ?
\end{enumerate}

\exo
Dans chaque cas, dire si le point appartient à la droite. Justifier.
\setlength{\columnseprule}{0pt}
\begin{multicols}{2}
\begin{enumerate}
\item $d_1: 2x-4y-10=0$ et $A(1;3)$.
\item $d_2: 3x+y -11 =0$ et $B(2;5)$.
\item $d_3: x=9$ et $C(9;2)$.
\end{enumerate}
\end{multicols}

\exo
Dans chaque cas, déterminer une équation cartésienne de la droite définie par le point et le vecteur directeur donnés.

\begin{enumerate}
\item $A(3;4)$ et  $\overrightarrow{u} \left( \begin{array}{c}
2\\
0
\end{array} \right)$

\item $B(-1;3)$ et  $\overrightarrow{v} \left( \begin{array}{c}
0\\
-1
\end{array} \right)$
\end{enumerate}

\exo
\begin{enumerate}
\item Déterminer une équation cartésienne de la droite $d$ passant par $A(-1;2)$ et de vecteur directeur $\overrightarrow{u} \left( \begin{array}{c}
2\\
3
\end{array} \right)$.
\item Tracer $d$.
\item Calculer les coordonnées des points d'intersection de $d$ avec les axes du repère, puis vérifier graphiquement.
\end{enumerate}

\exo
On donne les points $E(2;5)$ et $F(3;-1)$.
\begin{enumerate}
\item Tracer la droite $(EF)$, puis déterminer une équation cartésienne de cette droite.
\item Calculer l'ordonnée du point $K$, intersection de $(EF)$ et de l'axe des ordonnées.
\item Le point $M(-2;28)$ appartient-il à la droite $(EF)$ ?
\end{enumerate}

\exo
Pour chaque droite, donner un vecteur directeur, puis tracer la droite.
\begin{multicols}{2}
\begin{enumerate}
\item $4x+2y-1=0$
\item $-x+2y+2=0$
\item $2x+4=0$
\item $3y-2=0$
\end{enumerate}
\end{multicols}
\exo
Pour chacune des droites suivantes, déterminer son équation réduite.
\begin{multicols}{2}
\begin{enumerate}
\item $d_1: 2x+3y-1=0$
\item $d_2: 4x-2y+1=0$
\item $d_3:3y-6=0$
\item $d_4 : -3x+2y+1=0$
\end{enumerate}
\end{multicols}

\exo
Pour chacune des droites suivantes, déterminer un vecteur directeur, la pente lorsqu'elle existe, puis une équation cartésienne.

\begin{multicols}{2}
\begin{enumerate}
\item $d_1: y=4+\dfrac{1}{3}x$
\item $d_2: y=-\dfrac{3}{7}$
\item $d_3: y=\dfrac{7-4x}{5}$
\item $d_4:  x=-3$ 
\end{enumerate}
\end{multicols}

\exo
Dans chaque cas, déterminer si les droites $d$ et $d'$ sont sécantes ou parallèles.
Lorsqu'elles sont parallèles, préciser si elles sont confondues.

\begin{enumerate}
\item $d : 3x+4y+2=0$ et $d' : -6x-8y+3=0$
\item $d : 2x-3y+2=0$ et $d' : \dfrac{2}{5}x-\dfrac{3}{4} y+2=0$
\item $d : y=4x+2$ et $d' : \dfrac{1}{2}x-\dfrac{1}{8}y+\dfrac{1}{4}=0$
\end{enumerate}

\exo
Dans un parc zoologique, l'entrée coûte $30$ euros pour un adulte et $18$ euros pour un enfant.
À la fin de la journée, $630$ personnes ont visité le zoo et la recette est de $14220$ euros.
Parmi les personnes qui ont visité le zoo ce jour-là, combien y avait-il d'enfants ?
